\DocumentMetadata{
  pdfversion=2.0,
  pdfstandard=ua-2,
  lang=en-US,
  tagging=on,
  tagging-setup={math/setup=mathml-SE,tabsorder=structure}
}
\documentclass[11pt,letterpaper]{article}
\usepackage[margin=0.68in,includefoot]{geometry}
\usepackage{newcomputermodern}
\usepackage{amsmath,amscd,mathtools}
\usepackage{tikz,adjustbox}
\providecommand{\square}{\mdlgwhtsquare}
\usepackage[hidelinks]{hyperref}
\hypersetup{
  pdftitle={Algebra First-Year Exam, August 2025, Part 1},
  pdfauthor={University of Florida Department of Mathematics},
  pdfsubject={Graduate mathematics examination},
  pdfdisplaydoctitle=true,
  bookmarksopen=true
}
\setcounter{secnumdepth}{0}
\setlength{\parindent}{0pt}
\setlength{\parskip}{0.28em}
\setlength{\emergencystretch}{3em}
\setlength{\leftmargini}{2.15em}
\setlength{\leftmarginii}{2.65em}
\setlength{\leftmarginiii}{3em}
\pagestyle{empty}
\raggedbottom
\allowdisplaybreaks
\NewTaggingSocketPlug{block/list/label}{examproblem}{%
  \tagstructbegin{tag=Lbl}\tagstructend%
  \tagstructbegin{tag=\LBody}%
  \tagstructbegin{tag=\ProblemHeadingTag,title-o={\ProblemHeadingText}}%
  \tagmcbegin{tag=\ProblemHeadingTag,actualtext={\ProblemHeadingText}}%
  #2%
  \tagmcend%
  \tagstructend%
}
\newenvironment{examproblems}[2]{%
  \def\ProblemHeadingTag{#2}%
  \AssignTaggingSocketPlug{block/list/label}{examproblem}%
  \begin{enumerate}%
  \setcounter{enumi}{#1}%
  \renewcommand{\labelenumi}{\textbf{\theenumi.}}%
  \setlength{\topsep}{0.35em}%
  \setlength{\partopsep}{0pt}%
  \setlength{\itemsep}{0.48em}%
  \setlength{\parsep}{0.18em}%
}{\end{enumerate}}
\newenvironment{examsubparts}{%
  \AssignTaggingSocketPlug{block/list/label}{default}%
  \begin{enumerate}%
  \setlength{\topsep}{0.18em}%
  \setlength{\partopsep}{0pt}%
  \setlength{\itemsep}{0.16em}%
  \setlength{\parsep}{0.1em}%
}{\end{enumerate}}
\newenvironment{examinstructions}{%
  \par\begingroup\itshape
}{%
  \par\endgroup\vspace{0.18em}
}
\makeatletter
\renewcommand\subsection{\@startsection{subsection}{2}{0pt}%
  {-1.25ex plus -.25ex minus -.1ex}%
  {0.55ex plus .1ex}%
  {\normalfont\large\bfseries}}
\renewcommand{\@maketitle}{%
  \newpage\null\vskip -1em%
  {\footnotesize\bfseries
    \href{https://www.ufl.edu/}{University of Florida}, \href{https://math.ufl.edu/}{Department of Mathematics}\par}%
  \vskip -0.5em%
  \tagstructbegin{tag=H1,title={Algebra First-Year Exam, August 2025, Part 1}}%
  \tagpdfparaOff%
  \tagmcbegin{tag=H1}%
    {\LARGE\bfseries\href{https://gma.math.ufl.edu/exams/algebra-first-year/}{Algebra First-Year Exam}, August 2025, Part 1\par}%
  \tagmcend%
  \tagpdfparaOn%
  \tagstructend%
  \vskip 0.25em%
  \tagmcbegin{artifact}%
    \hrule height 1pt%
  \tagmcend%
  \vskip 1em%
}
\makeatother
\title{Algebra First-Year Exam, August 2025, Part 1}
\author{}
\date{}
\begin{document}
\maketitle
\thispagestyle{empty}
\begin{examinstructions}
Answer 4 questions. You should indicate which 4 problems you wish to be graded. Write your solutions clearly in complete English sentences. You may quote theorems (within reason) as long as you state them clearly.
\end{examinstructions}
\begin{examproblems}{0}{H2}
\def\ProblemHeadingText{Problem 1.}
\item
Prove that there exist at least three nonisomorphic nonabelian groups of order 12.
\def\ProblemHeadingText{Problem 2.}
\item
Prove that for \(n \geq 3\) the alternating group \(A_n\) is generated by 3-cycles. (You may assume that \(S_n\) is generated by transpositions.)
\def\ProblemHeadingText{Problem 3.}
\item
Let~\(p\) be a prime, and let \(G=\operatorname{GL}(3,\mathbb{F}_p)\), where \(\mathbb{F}_p\) denotes the field of~\(p\) elements.
\begin{examsubparts}
\item[\textnormal{(a)}]
Compute the order of~\(G\).
\item[\textnormal{(b)}]
Let
\[
U=\left\{\begin{pmatrix}1&a&b\\0&1&c\\0&0&1\end{pmatrix}\;\middle|\;a,b,c\in\mathbb{F}_p\right\}.
\]
 Prove that~\(U\) is a Sylow~\(p\)-subgroup of~\(G\).
\item[\textnormal{(c)}]
Find the commutator subgroup of~\(U\).
\end{examsubparts}
\def\ProblemHeadingText{Problem 4.}
\item
Let~\(G\) be a group acting on a set~\(X\).
\begin{examsubparts}
\item[\textnormal{(a)}]
For \(x\in X\), define the orbit of~\(x\) and the stabilizer of~\(x\).
\item[\textnormal{(b)}]
Assume that~\(G\) is finite. Prove that the size of the orbit of~\(x\) divides the order of~\(G\).
\end{examsubparts}
\def\ProblemHeadingText{Problem 5.}
\item
Let~\(G\) be a group.
\begin{examsubparts}
\item[\textnormal{(a)}]
Define the terms automorphism of~\(G\), inner automorphism of~\(G\) and characteristic subgroup of~\(G\).
\item[\textnormal{(b)}]
Prove that the set of automorphisms of~\(G\) is a group in which the subset of inner automorphisms forms a normal subgroup.
\item[\textnormal{(c)}]
Prove that if~\(N\) is a normal subgroup of~\(G\) and~\(H\) is a characteristic subgroup of~\(N\), then~\(H\) is normal in~\(G\).
\end{examsubparts}
\def\ProblemHeadingText{Problem 6.}
\item
Let~\(G\) be a group of order \(675=3^3 5^2\).
\begin{examsubparts}
\item[\textnormal{(a)}]
Prove that~\(G\) has a normal Sylow 5-subgroup.
\item[\textnormal{(b)}]
Prove that~\(G\) has an element of order 15.
\end{examsubparts}
\end{examproblems}
\end{document}
